\chapter{Quantum sensing of fundamental physics with superconducting circuits}
\label{ch:quantum-sensing}

\section*{Chapter overview}

Superconducting circuits combine three properties that make them
attractive single-quantum sensors of feeble signals from beyond the
Standard Model: macroscopic quantum coherence (a transmon flips
between $\ket{0}$ and $\ket{1}$ on absorbing one microwave quantum
of order $\hbar\omega_q\sim 20\,\mu\text{eV}$); gigahertz operating
frequencies that sit in the candidate-mass window for many
photon-coupled dark-matter species; and millikelvin operating
temperatures that suppress thermal backgrounds well below the
candidate signal energies. The same hardware that performs quantum
information processing therefore doubles as a particle-physics
front end, and a literature on cQED-based searches for dark matter
has emerged in the last decade.

This chapter is a guided tour of that literature. We begin with a
short primer on dark matter and on the kinematics of the two
recoil families (nuclear and electronic), specialise to the light,
photon-mediated candidates --- sub-GeV WIMP-like dark matter, dark
photons, and axions --- for which cQED is best suited, and then
walk through four sensing mechanisms based on superconducting
hardware:
\begin{itemize}
  \item phonon absorption mediated by quasiparticle creation
        (\cref{sec:phonon-qp}),
  \item infrared absorption in the Josephson junction
        (\cref{sec:ir-junction}),
  \item phonon-to-qubit coupling via piezoelectric substrates
        (\cref{sec:phonon-qubit-piezo}),
  \item direct microwave absorption in the qubit
        (\cref{sec:microwave-absorption}).
\end{itemize}
For each mechanism we sketch the signal model and quote a
sensitivity estimate. The chapter ends with the usual takeaways
box.

% --------------------------------------------------------------
\section{Dark matter, briefly}
\label{sec:dm-overview}

\paragraph{Why dark matter exists.}
Several independent observations --- galactic rotation curves, the
peaks of the cosmic-microwave-background power spectrum,
gravitational lensing of clusters, and the abundance of light
elements from big-bang nucleosynthesis --- agree on a universe
whose mass-energy budget is roughly 5\% baryonic matter, 27\%
non-luminous \emph{dark matter}, and 68\% dark energy. Dark matter
behaves gravitationally as cold, collisionless, non-relativistic,
non-baryonic matter; its particle physics, if any, is not in the
Standard Model and is the target of the searches we describe here.
A textbook account is given by Bertone, Hooper and
Silk~\cite{bertone2005}.

\paragraph{Local density and velocity.}
The Earth orbits within a roughly virialised dark-matter halo. The
locally relevant numbers are the mass density
\begin{equation}
\rho_\chi \;\approx\; 0.3\;\text{GeV}/\text{cm}^3,
\end{equation}
and a velocity distribution well approximated by a Maxwellian
truncated at the galactic escape speed, with one-dimensional
dispersion $\sigma_v\approx 156\,\text{km}/\text{s}\approx
5.2\times 10^{-4}\,c$ and a most-probable speed
$v_0\approx 220\,\text{km}/\text{s}\approx 7\times 10^{-4}\,c$. The
flux through a unit area on Earth is
\begin{equation}
\Phi_\chi \;\sim\; n_\chi\,\bar v
\;=\; \frac{\rho_\chi}{m_\chi}\,\bar v,
\end{equation}
which is large for light $\chi$ ($\sim 10^{12}\,\text{cm}^{-2}\,\text{s}^{-1}$
for $m_\chi=1\,\text{MeV}$) and small for heavy
($\sim 10^5\,\text{cm}^{-2}\,\text{s}^{-1}$ for $m_\chi=100\,\text{GeV}$).

\paragraph{Mass windows.}
The candidate mass spans nearly forty orders of magnitude. Three
broad bands matter for the present chapter:
\begin{itemize}
  \item \emph{Heavy WIMPs} ($m_\chi \sim 10\text{--}1000\,\text{GeV}$):
        targeted by traditional cryogenic phonon detectors and dual-phase
        time-projection chambers, where the signal is a nuclear or
        electronic recoil of $\mathcal{O}(\text{keV})$.
  \item \emph{Light WIMPs / sub-GeV} ($m_\chi \sim 10\,\text{keV}\text{--}1\,\text{GeV}$):
        recoil energies are at or below the eV scale; sub-eV
        thresholds, achievable with athermal-phonon and
        single-charge sensors built from superconducting hardware,
        are required.
  \item \emph{Ultralight bosons} ($m_\chi \sim 10^{-12}\text{--}10^{-3}\,\text{eV}$):
        wave-like dark matter coherently filling the lab volume,
        coupled to photons via either kinetic mixing (dark
        photons) or the QED anomaly (axions). The Compton
        frequency $f_\chi=m_\chi c^2/h$ falls naturally in the MHz
        to GHz range, matching superconducting resonators directly.
\end{itemize}

\paragraph{Why cQED?}
Superconducting circuits address the second and third bands with
the same toolset they already use for quantum information:
microwave-frequency operation, single-photon-level sensitivity,
millikelvin background, and quantum-limited amplification chains.
This chapter focuses on those bands.

% --------------------------------------------------------------
\section{Recoil signal models}
\label{sec:recoils}

\paragraph{Nuclear vs.\ electronic recoils.}
A dark-matter particle scattering elastically in a detector
deposits energy by recoiling against either a target nucleus
(nuclear recoil) or a target electron (electronic recoil). For an
incoming particle of mass $m_\chi$ and velocity $v_\chi$ on a
target of mass $M_T$, conservation of energy and momentum gives
the maximum recoil energy
\begin{equation}
\label{eq:recoil-max}
E_R^\text{max} \;=\; \frac{2\,\mu_{\chi T}^2\,v_\chi^2}{M_T},
\qquad
\mu_{\chi T} \equiv \frac{m_\chi M_T}{m_\chi+M_T},
\end{equation}
where $\mu_{\chi T}$ is the reduced mass. For a galactic
$v_\chi\sim 10^{-3}c$ and a silicon-target nucleus $M_T\approx 28\,m_p$,
$E_R^\text{max}$ is keV-scale only when $m_\chi\gtrsim$ a few GeV;
for $m_\chi\lesssim 100\,\text{MeV}$ it falls below an eV, and the
event becomes invisible to any threshold above that. Electronic
recoils are kinematically more favourable for light dark matter
because $M_T$ is the much smaller electron mass.

\paragraph{Threshold and rate.}
A detector with energy threshold $E_\text{th}$ and exposure
$\mathcal{E}=\text{(target mass)}\times\text{(live time)}$ is
sensitive to dark-matter particles whose differential rate
satisfies
\begin{equation}
\label{eq:rate}
R \;=\; \mathcal{E}\,\frac{n_\chi}{m_T}\,\langle\sigma_{\chi T}\,v\rangle,
\end{equation}
where $m_T$ is the target-particle mass, $\sigma_{\chi T}$ the
cross-section per target particle, and the brackets denote a
velocity average over the halo distribution. Lower $E_\text{th}$
opens up smaller $m_\chi$ via \cref{eq:recoil-max}; this is the
push that has driven the cQED programme. State-of-the-art
phonon-sensor thresholds reach $\mathcal{O}(\text{eV})$ today, with
a roadmap towards $\mathcal{O}(\text{meV})$ that exploits the
$\sim 0.34\,\text{meV}$ Cooper-pair binding energy of aluminium.

\paragraph{Light mediators and form factors.}
For a contact (heavy-mediator) interaction the velocity dependence
of $\sigma_{\chi T}$ is $\propto v^{-2}$ at low recoil energy. For
a light mediator (mass $\lesssim$ momentum transfer
$q\sim\mu_{\chi T}v$), an extra form factor $1/q^4$ enhances soft
recoils. The latter is what makes \emph{photon-mediated} dark
matter (next section) particularly soft and well-matched to a
microwave detector.

% --------------------------------------------------------------
\section{Photon-mediated dark matter}
\label{sec:photon-mediated-dm}

\paragraph{Dark photons.}
A dark photon $A'$ is a $U(1)'$ gauge boson kinetically mixed with
the Standard-Model photon~\cite{holdom1986},
\begin{equation}
\label{eq:dark-photon-mixing}
\mathcal{L} \;\supset\; -\tfrac{1}{4}F'_{\mu\nu}F'^{\mu\nu}
- \tfrac12 m_{A'}^2 A'_\mu A'^\mu
+ \tfrac{\varepsilon}{2}F'_{\mu\nu}F^{\mu\nu},
\end{equation}
with $\varepsilon$ the dimensionless mixing parameter. Field
redefinitions move the mixing into the matter sector, where the
dark photon couples to the electromagnetic current with effective
charge $\varepsilon e$. A non-zero $m_{A'}$ allows the dark photon
to be itself the dark matter; the local DM mass density then
corresponds to a coherent classical field
$\langle A'^2\rangle\propto\rho_\chi/m_{A'}^2$, oscillating at the
Compton frequency $f_{A'}=m_{A'}c^2/h$. For
$m_{A'}=4.1\,\mu\text{eV}$ this is $1\,\text{GHz}$, and a
microwave resonator looks the dark-matter field as a tiny coherent
RF source~\cite{caputo2021review}.

\paragraph{Axions.}
The axion is a pseudoscalar boson originally introduced to solve
the strong-$CP$ problem; it acquires a small mass and a
two-photon coupling
\begin{equation}
\label{eq:axion-coupling}
\mathcal{L}_{a\gamma\gamma} \;=\; -\tfrac14\,g_{a\gamma\gamma}\,a\,F_{\mu\nu}\tilde F^{\mu\nu}
\;=\; g_{a\gamma\gamma}\,a\,\bm E\!\cdot\!\bm B.
\end{equation}
In an external magnetic field $\bm B_0$, the axion field $a(t)$
sources an oscillating electric field via \cref{eq:axion-coupling}.
Inside a microwave cavity tuned to $\omega_a=m_a c^2/\hbar$, this
drives the resonant TM mode coherently, producing a tiny but
narrow-band signal --- the haloscope concept of
Sikivie~\cite{sikivie1983}.

\paragraph{Why microwave?}
Both $\rho_\chi/m_\chi^2\sim$ field amplitude and the mass--frequency
relation $f_\chi=m_\chi c^2/h$ make GHz the natural target band:
$1\,\text{GHz}\leftrightarrow 4.1\,\mu\text{eV}$, the QCD-axion
window of $\sim 1\text{--}40\,\mu\text{eV}$ corresponds to
$240\,\text{MHz}\text{--}10\,\text{GHz}$, and a transmon's
$\omega_q/2\pi\sim 5\,\text{GHz}$ sits inside it. Single-photon
sensitivity at GHz, achievable in cQED, gives a fundamental
advantage over heterodyne readout limited by quantum-noise
amplification.

% --------------------------------------------------------------
\section{Sensing mechanism 1: phonon--quasiparticle conversion}
\label{sec:phonon-qp}

\paragraph{Idea.}
A dark-matter particle scattering in a superconducting target
deposits energy as athermal phonons; a phonon of energy
$E_\text{ph}>2\Delta$ (where $2\Delta$ is the Cooper-pair binding
energy, $0.34\,\text{meV}$ in aluminium and $\sim 1.4\,\text{meV}$
in niobium) breaks a Cooper pair and creates two quasiparticles.
The quasiparticles diffuse to a transition-edge or microwave
kinetic-inductance sensor that reads them out as a small impedance
or temperature change~\cite{kurinsky2020}.

\paragraph{Threshold and reach.}
The sensitivity scales roughly as
\begin{equation}
E_\text{th} \;\sim\; 2\Delta\,\frac{1}{\eta_\text{collect}},
\end{equation}
with $\eta_\text{collect}$ the fraction of athermal phonons that
reach the sensor before downconverting to subgap energies. For Al
($\Delta\approx 170\,\mu\text{eV}$) and $\eta_\text{collect}\sim 0.1$,
$E_\text{th}\sim$ a few meV. This sets the kinematic floor at
$m_\chi\sim 10\,\text{MeV}$ for nuclear recoils (\cref{eq:recoil-max}
with $M_T=m_\text{Si nucleus}$, $v\sim 10^{-3}c$) --- already
significantly better than dual-phase TPC's $\sim$ GeV reach.

\paragraph{Sensitivity estimate.}
For a 1\,kg-day exposure with $E_\text{th}=10\,\text{meV}$ and
zero background, $\sigma_{\chi T}\sim 10^{-37}\,\text{cm}^2$ is
reachable for $m_\chi\sim 100\,\text{MeV}$. This is several orders
of magnitude beyond conventional WIMP searches at the same mass.

% --------------------------------------------------------------
\section{Sensing mechanism 2: infrared absorption in the junction}
\label{sec:ir-junction}

\paragraph{Idea.}
A Josephson junction itself is a Cooper-pair-breaking absorber: an
infrared or microwave photon of energy $\geq 2\Delta$, absorbed in
the leads, breaks a pair and injects a quasiparticle pair into the
junction, producing a measurable burst of poisoning. For a dark
photon of mass $m_{A'}\geq 2\Delta/c^2\approx 0.34\,\text{meV}/c^2$
(corresponding to $f_{A'}\geq 85\,\text{GHz}$), kinetic mixing
\cref{eq:dark-photon-mixing} converts the coherent dark-photon
field into ordinary photons inside the conductor, which then drive
quasiparticle creation~\cite{hochberg2021}.

\paragraph{Sensitivity to $\varepsilon$.}
The conversion rate scales as $\varepsilon^2$, the dark-photon
density as $\rho_\chi/m_{A'}^2$, and the absorption probability in
the lead as the product of absorber thickness and the Mattis--Bardeen
conductivity. Order-of-magnitude, a 1\,m$^2$-cm-thick aluminium
absorber operated for 1\,year reaches
$\varepsilon\sim 10^{-13}$ at $m_{A'}\sim 1\,\text{meV}$, an order
of magnitude beyond solar lifetime constraints in that mass band.
The shape of the reach curve $\varepsilon(m_{A'})$ is set by the
absorber's frequency-dependent surface impedance.

% --------------------------------------------------------------
\section{Sensing mechanism 3: phonon--qubit coupling via piezoelectricity}
\label{sec:phonon-qubit-piezo}

\paragraph{Idea.}
A phonon in a piezoelectric substrate (AlN, GaAs, sapphire-based
sandwich structures) carries an electric polarisation that couples
linearly to a transmon's charge, $\op{n}$. Single phonons of energy
in the 1--10\,GHz band can therefore drive transmon transitions
in exact analogy with single microwave photons in cavity QED. The
qubit thus acts as a single-phonon detector, with the phonon
playing the role of the readout
quantum~\cite{karasik2014,linehan2024}.

\paragraph{Coupling strength and detection.}
The relevant Hamiltonian is, schematically,
\begin{equation}
\op{H}_\text{int} \;=\; \hbar g\,\hat\sigma_x\,(\hat b + \hat b^\dagger),
\end{equation}
with $\hat b$ the phonon-mode annihilation operator and
$g\propto e_{14}\,d/\sqrt{\hbar/(2 M\Omega_\text{ph})}$ depending
on the piezoelectric coefficient $e_{14}$, the qubit-phonon overlap
length $d$, and the phonon mode mass and frequency. Typical
$g/2\pi\sim 0.1\text{--}10\,\text{MHz}$. Detecting a single phonon
with a transmon at $\omega_q/2\pi=5\,\text{GHz}$ is therefore
competitive with conventional KID-based detection above the
sub-meV phonon-energy threshold.

\paragraph{Sensitivity reach.}
For sub-GeV nuclear-recoil dark matter on Si, a piezoelectric
target plus transmon readout reaches
$\sigma_{\chi T}\sim 10^{-38}\,\text{cm}^2$ at
$m_\chi\sim 1\,\text{MeV}$ for a few kg-day exposure --- the lowest
mass reach yet projected for nuclear recoils.

% --------------------------------------------------------------
\section{Sensing mechanism 4: microwave absorption in the qubit}
\label{sec:microwave-absorption}

\paragraph{Idea.}
A transmon driven by an oscillating microwave field at its
transition frequency $\omega_q$ has an absorption cross-section set
by its dipole moment. If the field is sourced by a dark-matter
candidate --- a kinetically mixed dark photon, or an axion in a
strong static $\bm B_0$ --- a population transfer
$\ket{0}\to\ket{1}$ counts as a candidate event. The transmon
operates as a single-microwave-photon counter in the 1--20\,GHz
window~\cite{dixit2021}.

\paragraph{Sensitivity formula.}
For a dark photon of mass $m_{A'}$ and a cavity of volume $V$
hosting a TM mode with form factor $C\sim O(1)$, the equivalent
ordinary-photon power on resonance is
\begin{equation}
P_\text{sig} \;=\; \tfrac{1}{4}\,\varepsilon^2\,\rho_\chi\,V\,C\,\omega_{A'}.
\end{equation}
A transmon with photon-counting efficiency $\eta\sim 0.5$ and
dark-count rate $\Gamma_\text{dc}\sim 1\,\text{Hz}$, integrating
for $\tau$, reaches the SNR-1 limit
\begin{equation}
\varepsilon^2_\text{min}
\;\sim\; \frac{\sqrt{\Gamma_\text{dc}/\tau}}{\eta\,(\rho_\chi V C\omega_{A'})}.
\end{equation}
Plugging the numbers: 1\,litre cavity, $\omega_{A'}/2\pi=5\,\text{GHz}$,
$\tau=10$\,days, $\eta=0.5$, $\Gamma_\text{dc}=1\,\text{mHz}$:
$\varepsilon_\text{min}\sim 10^{-15}$, which beats heterodyne
amplification by $1\text{--}2$ orders of magnitude in this band.
Repeating across the cavity tuning range fills out the
$\varepsilon(m_{A'})$ reach.

\paragraph{Axion case.}
For axions in a static $\bm B_0$, the cavity power is the standard
Sikivie haloscope formula
$P_\text{sig}\propto g_{a\gamma\gamma}^2\,\rho_a\,B_0^2\,V\,Q_c$,
which a single-photon-counting transmon converts into a count rate.
The same accounting reaches $g_{a\gamma\gamma}\sim 10^{-15}\,\text{GeV}^{-1}$
for a 1\,T field over a few-litre cavity, sweeping the QCD axion
window.

% --------------------------------------------------------------
\section*{Chapter takeaways and outlook}

\begin{chaptersummary}

\paragraph{Why cQED for fundamental physics.}
Macroscopic quantum coherence at GHz frequencies, single-microwave-photon
sensitivity, and millikelvin operating temperature put
superconducting hardware in a unique position to detect feeble
signals from beyond the Standard Model in the sub-GeV and
ultralight-boson dark-matter mass bands.

\paragraph{Recoil kinematics.}
For elastic dark-matter scattering on a target of mass $M_T$, the
maximum recoil energy is
$E_R^\text{max}=2\mu_{\chi T}^2 v_\chi^2/M_T$
(\cref{eq:recoil-max}). At galactic velocities $v_\chi\sim 10^{-3}c$
this drops below an eV for $m_\chi\lesssim 100\,\text{MeV}$,
demanding sub-eV thresholds.

\paragraph{Photon-mediated dark matter.}
Dark photons couple to the electromagnetic current with effective
charge $\varepsilon e$ via the kinetic-mixing Lagrangian
\cref{eq:dark-photon-mixing}. Axions couple to two photons with
strength $g_{a\gamma\gamma}$ \cref{eq:axion-coupling}, in a
constant $\bm B_0$ producing an oscillating $\bm E$ at the Compton
frequency $f_a=m_a c^2/h$. The mass--frequency relation puts the
QCD-axion window in the 240\,MHz--10\,GHz range, exactly where
transmons live.

\paragraph{Four sensing mechanisms.}
\begin{itemize}
  \item Phonon-quasiparticle conversion (\cref{sec:phonon-qp}) reaches
        $E_\text{th}\sim$ meV via Cooper-pair breaking
        ($2\Delta\approx 0.34\,\text{meV}$ in Al), opening
        $m_\chi\sim 10\,\text{MeV}$.
  \item Infrared absorption in the junction
        (\cref{sec:ir-junction}) reaches dark-photon mixing
        $\varepsilon\sim 10^{-13}$ at $m_{A'}\sim 1\,\text{meV}$
        with kg-scale Al absorbers.
  \item Phonon-qubit coupling via piezoelectrics
        (\cref{sec:phonon-qubit-piezo}) uses transmons as
        single-phonon detectors with $g/2\pi\sim 1\,\text{MHz}$,
        targeting nuclear recoils down to $m_\chi\sim 1\,\text{MeV}$.
  \item Microwave absorption in the qubit
        (\cref{sec:microwave-absorption}) implements a
        single-photon-counting haloscope, reaching
        $\varepsilon\sim 10^{-15}$ for dark photons and the QCD-axion
        window for $g_{a\gamma\gamma}$.
\end{itemize}

\paragraph{Open questions.}
The numbers above are best-case projections; the limiting
backgrounds in practice are quasiparticle poisoning, two-level-system
loss in the substrate, and stray IR/microwave radiation through the
cryostat. Reducing each is an active programme, on the same critical
path as scaling up quantum-information cQED itself.

\end{chaptersummary}
